
\documentclass[a4paper,12pt]{article}
\usepackage[utf8]{inputenc}
\usepackage[T1]{fontenc}
\usepackage[french]{babel}
\usepackage{geometry}
\geometry{margin=2cm}
\usepackage{setspace}
\onehalfspacing

\title{Étude de Risques - Société Educ\&Info}
\author{DSI - BTS SIO}
\date{}

\begin{document}

\maketitle

\section*{1. Identification des processus métiers / actifs informationnels}
\begin{itemize}
    \item Conception et mise à jour des programmes de formation.
    \item Gestion administrative : facturation, paiements, relations avec les OPCA.
    \item Gestion commerciale : devis, suivi client, communication avec le site web.
    \item Organisation logistique des formations : planification, préparation des salles, matériel.
    \item Gestion du site internet : publication des cours, inscriptions en ligne.
    \item Gestion des supports pédagogiques et évaluations.
    \item Suivi post-formation et adaptation des contenus selon les retours clients.
\end{itemize}

\section*{2. Identification des actifs supports}
\begin{itemize}
    \item Infrastructure réseau : serveurs internes (RADIUS, base de données, sauvegarde), routeur, pare-feu.
    \item Postes de travail du personnel et des stagiaires.
    \item ERP pour la gestion des plannings, facturation et suivi des formations.
    \item Site internet hébergé sur un serveur externe.
    \item Données clients, documents pédagogiques, données administratives.
    \item Salle informatique sécurisée et climatisée.
    \item Connexion VPN avec la société mère.
\end{itemize}

\section*{3. Identification des mesures de sécurité existantes}
\begin{itemize}
    \item Moyens anti-incendie présents dans les locaux.
    \item Alarme anti-intrusion active en dehors des heures d’ouverture.
    \item Accès au réseau Wi-Fi via portail captif et serveur RADIUS.
    \item Sauvegarde quotidienne des fichiers sur un serveur interne.
    \item Transmission hebdomadaire des fichiers administratifs via VPN.
    \item Armoires métalliques fermant à clé pour les documents papier.
    \item Maintenance logicielle assurée sous 24h par contrat fournisseur.
\end{itemize}

\section*{4. Identification des vulnérabilités}
\begin{itemize}
    \item Absence de PSSI et de charte d’utilisation du SI.
    \item Aucun contrôle d’accès sur les postes stagiaires.
    \item Absence de procédure de fermeture et de contrôle des clés des salles.
    \item Utilisation de matériels personnels connectés au réseau de l’entreprise.
    \item Réseau saturé lorsque toutes les salles sont utilisées.
    \item Manque de sensibilisation à la sécurité du personnel.
    \item DSI non expert en sécurité informatique.
\end{itemize}

\section*{5. Identification des menaces}
\begin{itemize}
    \item Intrusion physique ou vol de matériel informatique.
    \item Intrusion réseau ou vol de données (clients, supports pédagogiques).
    \item Infection par malware due à l’utilisation de périphériques externes.
    \item Perte de données liée à une défaillance matérielle ou erreur humaine.
    \item Déni de service lié à la saturation du réseau.
    \item Espionnage industriel ou copie non autorisée de supports de cours.
    \item Risques liés à l’intervention de prestataires externes (nettoyage, maintenance).
\end{itemize}

\section*{6. Identification des impacts}
\begin{itemize}
    \item Perte ou corruption de données critiques ou confidentielles.
    \item Interruption de l’activité de formation.
    \item Atteinte à la réputation de l’entreprise.
    \item Pertes financières liées à des incidents ou litiges clients.
    \item Perte de confiance des clients et partenaires.
    \item Non-conformité réglementaire (protection des données).
\end{itemize}

\section*{7. Identification des risques}
\begin{itemize}
    \item Risque de fuite ou de vol de données sensibles (clients, supports de cours).
    \item Risque d’intrusion sur le réseau stagiaire ou personnel.
    \item Risque de panne matérielle entraînant la perte de disponibilité du SI.
    \item Risque de dégradation de l’image et perte de clients.
    \item Risque de non-conformité en matière de protection des données.
    \item Risque d’exploitation de vulnérabilités dues à l’absence de PSSI.
\end{itemize}

\end{document}
